Свяжем решения систем $Ax=b$ (1) и $Ax=0$ (2).
\begin{proposition}
    Пусть $x_0$ - решение системы (1). $\tilde{x}$ является решением системы (1) тогда и только тогда, когда существует $y$ - решение системы (2), для которого выполнено $\tilde{x} = x_0 +y$.
\end{proposition}
\begin{proof}
    Докажем утверждение в обе стороны отдельно:
    \begin{itemize}
        \item[$\Rightarrow$] Если $\tilde{x}$ - решение системы (1), то имеем $Ax_0=b$ и $A\tilde{x}=b$, следовательно $A(\tilde{x} - x_0)=0$. Значит в качестве $y$ можем взять $\tilde{x}-x_0$, откуда $\tilde{x}=x_0 + y$.
        \item[$\Leftarrow$] Если $y$ - решение системы (2), то имеем $A\tilde{x}=A(x_0+y)=Ax_0+Ay=b+0=b$. Значит $\tilde{x}$ - решение системы (1).
    \end{itemize}
\end{proof}
\textit{Вывод}: если мы умеем решать систему (2) и имеем хотя бы одно решение системы (1), то мы получаем все решения системы (1).
\begin{proposition}
    Пусть $h_1, h_2$ - решения системы $Ax=0$, тогда $\forall \alpha, \beta \in \R \hookrightarrow h=\alpha h_1 + \beta h_2$ - тоже решение.
\end{proposition}
\begin{proof}
    Из $A(\alpha h_1 + \beta h_2)=\alpha A h_1 + \beta A h_2=0$ получаем требуемое.
\end{proof}
\begin{definition}
\label{Фундаментальная система решений.}
    Матрица $F$, состоящая из столбцов высоты $n$, называется фундаментальной матрицей системы $Ax=0$, если выполняются три условия:
    \begin{enumerate}
        \item $AF=0$ (т.е. каждый из столбцов $F$ - решение системы $Ax=0$).
        \item Столбцы $F$ - ЛНЗ.
        \item $rg(F)$ максимален среди всех матриц, удовлетворяющих условию (1).
    \end{enumerate}
\end{definition}
Идея в том, что любое решение системы выражается линейной комбинацией столбцов матрицы $F$, так как иначе ранг $F$ не максимален (к найденной $F$ можем приписать столбец, не являющийся линейной комбинацией ее столбцов, увеличив ранг).
Если $rg(A)=n$, то из ЛНЗ столбцов матрицы $A$ получаем, что у однородной системы только нулевое решение и матрица $F$ не существует.
\begin{theorem}
    Пусть дана система $Ax=0$ и её фундаментальная матрица $F$, тогда \[rg(F) \leqslant n-rg(A)\]
\end{theorem}
\begin{proof} Приведем матрицу $A$ элементарными операциями со строками и перестановкой столбцов к простейшему виду. Зануляя небазисные столбцы получим следующий вид:
\[A \longrightarrow A'=
\left(
    \begin{array}{c|c}
    E & 0 \\
    \hline
    0 & 0
    \end{array}
\right)
\]
Пусть $P, Q$ - произведения элементарных матриц, соответствующих нашим преобразованиям матрицы $A$. Тогда $A'=PAQ$. Рассмотрим вспомогательную матрицу $F'=Q^{-1}F$. Так как $Q$ невырожденная, то по утверждению 2.4 $rg(F')=rg(F)$. 

Далее несложно заметить, что у $A'F'$ первые $rg(A)$ строк совпадают с первыми $rg(A)$ строками $F'$. Но $A'F'=PAQQ^{-1}F=P(AF)=0$. Значит первые $rg(A)$ строк $F'$ являются нулевыми, а значит $rg(F)=rg(F') \leqslant n-rg(A)$.
\end{proof}
Поймем, что подставляя в систему $A'x=0$ столбцы, у которых последние $n-rg(A)$ переменных фиксированные и являющиеся набором из $n-rg(A) -1$ нуля и одной единицы, мы получаем ровно $n-rg(A)$ ЛНЗ решений системы. А значит эти решения образуют фундаментальную матрицу $F$ с рангом $n-rg(A)$.
\[
\begin{array}{cccc}
\begin{pmatrix}
\vdots \\
1 \\
0 \\
\vdots \\
0
\end{pmatrix} & 
\begin{pmatrix}
\vdots \\
0 \\
1 \\
\vdots \\
0
\end{pmatrix} & 
\cdots & 
\begin{pmatrix}
\vdots \\
0 \\
\vdots \\
0 \\
1
\end{pmatrix}
\end{array}
\hookrightarrow F
\]
\begin{proposition}
\label{Общее решение однородной системы.}
    Пусть дана система $Ax=0$, тогда 
    \begin{itemize}
        \item для любого столбца $c$ высоты $n-r$,  $Fc$ - решение.
        \item для любого решения $\tilde{x}$ существует столбец $c$ такой, что $Fc=\tilde{x}$.
    \end{itemize}
\end{proposition}
\begin{proof}
    Первая часть утверждения очевидно, так как $A(Fc)=(AF)c=0c=0$. Во второй части: так как по доказанному ранее $\tilde{x}$ представляется в виде линейной комбинации столбцов $F$, то вектор из коэффициентов этой комбинации можно взять в качестве $c$.
\end{proof}
Резюмируем:
\begin{proposition}
\label{Общее решение неоднородной системы.}
    Если даны система $Ax=b$ (1), её частное решение $x_0$, фундаментальная матрица $F$ для системы $Ax=0$ (2), то общее решение системы (1) есть сумма её частного решения и общего решения системы (2) \[x=x_0+Fc\]
\end{proposition}
\section{Линейное пространство}
\begin{definition}
\label{Аксиоматика линейного пространства.}
    Множество $L$ будем называть линейным пространством, а его элементы векторами, если на нём заданы операция сложения $\forall x,y \in L \hookrightarrow x+y \in L$ и операция умножения на число $\forall x \in L, \alpha \in \R \hookrightarrow \alpha x \in L$, и выполняются следующие аксиомы:
    \begin{enumerate}
        \item $\forall x, y \in L \hookrightarrow x+y=y+x$
        \item $\forall x, y, z \in L \hookrightarrow (x+y)+z=x+(y+z)$
        \item $\exists 0\in L:\forall x\in L \hookrightarrow x+0=x$
        \item $\forall x \in L:\exists (-x) \in L \hookrightarrow x+(-x)=0$
        \item $\forall \alpha, \beta \in \R: \forall x \in L \hookrightarrow (\alpha \beta)x=\alpha (\beta x)$
        \item $\forall x \in L \hookrightarrow 1x=x$
        \item $\forall \alpha, \beta \in \R: \forall x \in L \hookrightarrow (\alpha+\beta)x=\alpha x+\beta x$
        \item $\forall \alpha \in \R: \forall x, y \in L \hookrightarrow \alpha(x+y)=\alpha x+\alpha y$
    \end{enumerate}
\end{definition}
\textit{Комментарий}: из определения можно сделать следующие наблюдения
\begin{itemize}
    \item Множество $L$ непустое по аксиоме 3.
    \item Из аксиом 1, 2, 3, 4 линейное пространство - Абелева группа по сложению.
    \item Множество решений однородной системы всегда образуют линейное пространство.
    \item Единственность 0 и противоположного элемента - следствия из аксиом.
\end{itemize}

Понятия ЛЗ и ЛНЗ векторов вводятся аналогично нашим предыдущим определениям. Все теоремы о ЛЗ/ЛНЗ векторов остаются в силе.